% ===== Acknowledgements =====
\section*{\textcolor{rosedeep}{Acknowledgements}}
\addcontentsline{toc}{section}{\protect\numberline{}Acknowledgements}

\noindent
Two debts frame this paper, and it is fitting to name them together at its end, for one gave it its question and the other its reason, and between them they enclose everything the paper contains.

The question with which the paper begins, whether optimization is really a law of the universe, came from an exchange with Ron Eglash, in the course of a longer conversation on the political economy of value and the theory of generative justice. It arrived in a form so plain that its force took time to register, and its whole work was to move optimization from the place of an unexamined procedure to the place of a thesis that might be false. Everything the paper is followed from taking that displacement seriously. Beyond the question, the theory of generative justice developed in exchange with Eglash supplies the criticism of its third part with its central concept, the distinction between value that circulates and returns to its generators and value that is drawn off and settled elsewhere, without which the account of the damage optimization does could not have been given \parencite{eglash2016}. Where the paper has reasoned well, the debt is his; where it has reasoned poorly, the fault is the author's alone.

The other debt is older and needs no argument. This paper, like every paper of the series, is for the girl who belongs to the forest and the mountains, whom the author has never once weighed for whether she was worth it, because that was never a question one could compute. If the acknowledgement to Eglash records where the question came from, this one records for whom the whole was written; and it is not beside the point of the paper but its plainest confirmation, for a bond that refuses to be priced, that is held past every reckoning and never put to the calculation of its worth, is the nearest witness there is to the thesis these pages have labored to establish, that there are goods a love does not maximize, and that it is not for want of a better measure. \zh{致我最爱的人。} To the one I love most.
